\section{模型假设与符号说明}

\subsection{模型假设}

% 正文采用编号式论证，不把假设默认排成三线表。类型、可验证性、失败影响等完整审计字段
% 仍保存在 reports/methodology/statistical_assumptions.json，写作时只把与评委阅读有关的
% “具体假设 + 题目依据 + 建模作用”转换为自然语言。
%
% 由生成器按题目实际内容写入 3--6 条，每条为一个自然段。例如调用：
% \assumptionnote{一（短名称）}{具体、可检验或有明确边界的假设。}{题面依据以及该假设在模型中的作用。}
% 最终显示为“假设一（短名称）：具体假设。依据与作用自然接续。”，不另起审计字段行。
%
% 禁止预置“数据真实可靠”“忽略所有交互作用”“变量连续可微”“缺失值随机缺失”等
% 万能假设；evidence_type=untestable_from_provided_data 时禁止写“经验证成立”。
% 对结论影响重大的失效情形，放入灵敏度分析或模型评价，不在此处机械堆成五列表格。

% 由 3_assumptions.tex 在第三章内调用；保留独立片段以兼容旧项目入口。
\subsection{符号说明}

% v2（§13.5 SymbolTable）：符号/含义/单位三列语义固定、longtable 可跨页（T147）；
% 禁止整张符号表因不能分页被整体推到下一页造成前页半空。
\SymbolTable{主要符号说明}{
  % TODO: 从 FINAL_MODEL_SPEC 的 variables/symbols 生成行（模板禁止预置 N/M/X/y/R² 示例符号）。
  % 行格式：$N$ & 样本数量 & 个 \\
}

